
\documentclass[11pt,a4paper]{article}
\usepackage{Act}
\begin{document}
\input{\detokenize{/home/fenarius/Travail/Cours/cpge-info/latex/Macros.tex}}
\RDS{4}{14/06/2026}

\begin{tcolorbox}[enhanced,width=\textwidth,center upper,fontupper=\bfseries,drop shadow southwest,sharp corners]
{\sc \large Astruc} Alexandre
\end{tcolorbox}
\medskip
\begin{tabularx}{\textwidth}{p{5cm}X}
	\alertbox{\faAward}{Note}{
		\begin{itemize}[leftmargin=0pt]
			\item[\textbullet] Note : \textbf{\large 11.6}
			\item[\textbullet] Rang : \textbf{8}
			\item[\textbullet] Traité : 83 \%
		\end{itemize}
	} &
	\alertbox{\faChartLine}{Statistiques des notes}{
        \psset{xunit=1cm, yunit=1cm,fillstyle=solid}
		\begin{pspicture}(0,-0.1)(16,1.45)
		   \savedata{\data}[11.6 14.5 18.8 8.4 8.2 7.6 14.1 8.1 9.6 14.2 17.6 14.9 18.6 6.4 0.0]
		   \rput{-90}(0,0.9){\psBoxplot[barwidth=1.1cm,yunit=0.5,fillcolor=gray,linewidth=1pt]{\data}}
		   \psaxes[yAxis=false,dx=1cm,Dx=2,labelsep=1pt,linecolor=gray,xlabelFontSize=\scriptstyle](0,0)(10.1,4)
		   \psdot[dotsize=8pt,dotstyle=diamond,linecolor=black,fillstyle=solid,fillcolor=white,linewidth=1pt](5.8,0.85)
           \psdot[dotsize=6pt,dotstyle=x,linecolor=black,linewidth=3pt](5.753333333333333,0.85)
		   \end{pspicture}
	} \\
    
\end{tabularx}\\
\begin{tabularx}{\textwidth}{X}
\alertbox{\faComment}{Commentaire}
{
	Bon travail, tu es en progrès il faut continuer sur cette voie, c’est très bien !
Sur certaines questions tu sembles avoir les bonnes idées mais la rédaction de tes justifications sont encore à travailler !
}
\end{tabularx}
\medskip
    \ding{113} \textbf{\sffamily{Résultats par thème}} \medskip \\
    \renewcommand{\arraystretch}{1.2}
    \begin{tabular}{|l|r|r|}
    \cline{2-3}
    \multicolumn{1}{l|}{} & \multicolumn{1}{|c|}{Points} & \multicolumn{1}{|c|}{Traitées} \\
    \hline
    {Base d'OCaml (fonctionnel)} & 24\% \;{\small (11/45)} & 60\% \;{\small (3/5)} \\ \hline {Complexité d'un algorithme} & 46\% \;{\small (14/30)} & 75\% \;{\small (3/4)} \\ \hline {Représentation des entiers} & 100\% \;{\small (05/5)} & 100\% \;{\small (1/1)} \\ \hline {Structure de données séquentielles} & 82\% \;{\small (37/45)} & 100\% \;{\small (6/6)} \\ \hline {Base de données et SQL} & 53\% \;{\small (32/60)} & 100\% \;{\small (8/8)} \\ \hline {Comprendre un algorithme} & 75\% \;{\small (15/20)} & 75\% \;{\small (3/4)} \\ \hline {Types structurés en C et pointeurs} & 100\% \;{\small (20/20)} & 100\% \;{\small (2/2)} \\ \hline {Table de hachage} & 40\% \;{\small (16/40)} & 50\% \;{\small (2/4)} \\ \hline {Programmation de base en C} & 86\% \;{\small (13/15)} & 100\% \;{\small (2/2)} \\ \hline \end{tabular} \\\\\medskip \\
    \ding{113} \textbf{\sffamily{Résultats par exercice}} \medskip \\
    \renewcommand{\arraystretch}{1.2}
    \begin{tabular}{|l|r|r|}
    \cline{2-3}
    \multicolumn{1}{l|}{} & \multicolumn{1}{|c|}{Points} & \multicolumn{1}{|c|}{Traitées} \\
    \hline
    Exercice {1} & 70\% \;{\small (39/55)} & 100\% \;{\small (7/7)} \\ \hline Exercice {2} & 76\% \;{\small (38/50)} & 85\% \;{\small (6/7)} \\ \hline Exercice {3} & 53\% \;{\small (32/60)} & 100\% \;{\small (8/8)} \\ \hline Exercice {4} & 56\% \;{\small (31/55)} & 66\% \;{\small (4/6)} \\ \hline Exercice {5} & 38\% \;{\small (23/60)} & 62\% \;{\small (5/8)} \\ \hline \end{tabular} \\\\
    \vspace{0.5cm}\\
    \ding{113} \textbf{\sffamily{Historique des notes}} \medskip \\
    \psset{xunit=1.4cm, yunit=0.2cm}
    \begin{pspicture}(-1,-1)(12,22)

% --- Axe X : numéros de devoir ---
\psaxes[
    Dx=1,
    Dy=2,
    Ox=0,
    Oy=0,
    labels=all,
    ticks=all,
](0,0)(0,0)(7,20)


\listplot[plotstyle=line,showpoints=true,linecolor=black,linewidth=0.7pt, dotstyle=diamond*, dotsize=0.2]{%
    1 11.7
2 11.9
3 10.3
4 11.6
}

% Minimum de la classe
\listplot[plotstyle=line,showpoints=true,linecolor=gray,linewidth=0.7pt, dotstyle=|, dotangle=90, dotsize=0.15, linestyle=dotted]{%
    1 5.8
2 4.6
3 0.9
4 0
}

% Maximum de la classe
\listplot[plotstyle=line,showpoints=true,linecolor=gray,linewidth=0.7pt, dotstyle=|, dotangle=90, dotsize=0.15, linestyle=dotted]{%
    1 18.9
2 18.4
3 19.7
4 18.8
}

% Moyenne de la classe
\listplot[plotstyle=line,showpoints=true,linecolor=gray,linewidth=0.7pt, dotstyle=x, dotsize = 0.15, linestyle=dashed]{%
    1 12.4466666666667
2 13.1933333333333
3 10.82
4 11.5066666666667
}
\end{pspicture}
\pagebreak
\begin{tcolorbox}[enhanced,width=\textwidth,center upper,fontupper=\bfseries,drop shadow southwest,sharp corners]
{\sc \large Berfeuil} Rohan
\end{tcolorbox}
\medskip
\begin{tabularx}{\textwidth}{p{5cm}X}
	\alertbox{\faAward}{Note}{
		\begin{itemize}[leftmargin=0pt]
			\item[\textbullet] Note : \textbf{\large 14.5}
			\item[\textbullet] Rang : \textbf{5}
			\item[\textbullet] Traité : 92 \%
		\end{itemize}
	} &
	\alertbox{\faChartLine}{Statistiques des notes}{
        \psset{xunit=1cm, yunit=1cm,fillstyle=solid}
		\begin{pspicture}(0,-0.1)(16,1.45)
		   \savedata{\data}[11.6 14.5 18.8 8.4 8.2 7.6 14.1 8.1 9.6 14.2 17.6 14.9 18.6 6.4 0.0]
		   \rput{-90}(0,0.9){\psBoxplot[barwidth=1.1cm,yunit=0.5,fillcolor=gray,linewidth=1pt]{\data}}
		   \psaxes[yAxis=false,dx=1cm,Dx=2,labelsep=1pt,linecolor=gray,xlabelFontSize=\scriptstyle](0,0)(10.1,4)
		   \psdot[dotsize=8pt,dotstyle=diamond,linecolor=black,fillstyle=solid,fillcolor=white,linewidth=1pt](7.25,0.85)
           \psdot[dotsize=6pt,dotstyle=x,linecolor=black,linewidth=3pt](5.753333333333333,0.85)
		   \end{pspicture}
	} \\
    
\end{tabularx}\\
\begin{tabularx}{\textwidth}{X}
\alertbox{\faComment}{Commentaire}
{
	C’est  bien, le SQL notamment a l’air bien maitrisé. Attention cependant à la définition d’un arbre binaire de recherche. La propriété doit être vraie pour TOUS les nœuds et elle indique que TOUTES les valeurs présentes à gauche (resp. droite) sont plus petites (resp. grandes).
}
\end{tabularx}
\medskip
    \ding{113} \textbf{\sffamily{Résultats par thème}} \medskip \\
    \renewcommand{\arraystretch}{1.2}
    \begin{tabular}{|l|r|r|}
    \cline{2-3}
    \multicolumn{1}{l|}{} & \multicolumn{1}{|c|}{Points} & \multicolumn{1}{|c|}{Traitées} \\
    \hline
    {Base d'OCaml (fonctionnel)} & 66\% \;{\small (30/45)} & 100\% \;{\small (5/5)} \\ \hline {Complexité d'un algorithme} & 53\% \;{\small (16/30)} & 75\% \;{\small (3/4)} \\ \hline {Représentation des entiers} & 20\% \;{\small (01/5)} & 100\% \;{\small (1/1)} \\ \hline {Structure de données séquentielles} & 68\% \;{\small (31/45)} & 100\% \;{\small (6/6)} \\ \hline {Base de données et SQL} & 100\% \;{\small (60/60)} & 100\% \;{\small (8/8)} \\ \hline {Comprendre un algorithme} & 65\% \;{\small (13/20)} & 100\% \;{\small (4/4)} \\ \hline {Types structurés en C et pointeurs} & 30\% \;{\small (06/20)} & 50\% \;{\small (1/2)} \\ \hline {Table de hachage} & 90\% \;{\small (36/40)} & 100\% \;{\small (4/4)} \\ \hline {Programmation de base en C} & 66\% \;{\small (10/15)} & 50\% \;{\small (1/2)} \\ \hline \end{tabular} \\\\\medskip \\
    \ding{113} \textbf{\sffamily{Résultats par exercice}} \medskip \\
    \renewcommand{\arraystretch}{1.2}
    \begin{tabular}{|l|r|r|}
    \cline{2-3}
    \multicolumn{1}{l|}{} & \multicolumn{1}{|c|}{Points} & \multicolumn{1}{|c|}{Traitées} \\
    \hline
    Exercice {1} & 63\% \;{\small (35/55)} & 100\% \;{\small (7/7)} \\ \hline Exercice {2} & 42\% \;{\small (21/50)} & 57\% \;{\small (4/7)} \\ \hline Exercice {3} & 100\% \;{\small (60/60)} & 100\% \;{\small (8/8)} \\ \hline Exercice {4} & 85\% \;{\small (47/55)} & 100\% \;{\small (6/6)} \\ \hline Exercice {5} & 66\% \;{\small (40/60)} & 100\% \;{\small (8/8)} \\ \hline \end{tabular} \\\\
    \vspace{0.5cm}\\
    \ding{113} \textbf{\sffamily{Historique des notes}} \medskip \\
    \psset{xunit=1.4cm, yunit=0.2cm}
    \begin{pspicture}(-1,-1)(12,22)

% --- Axe X : numéros de devoir ---
\psaxes[
    Dx=1,
    Dy=2,
    Ox=0,
    Oy=0,
    labels=all,
    ticks=all,
](0,0)(0,0)(7,20)


\listplot[plotstyle=line,showpoints=true,linecolor=black,linewidth=0.7pt, dotstyle=diamond*, dotsize=0.2]{%
    1 11.2
2 15.4
3 16.4
4 14.5
}

% Minimum de la classe
\listplot[plotstyle=line,showpoints=true,linecolor=gray,linewidth=0.7pt, dotstyle=|, dotangle=90, dotsize=0.15, linestyle=dotted]{%
    1 5.8
2 4.6
3 0.9
4 0
}

% Maximum de la classe
\listplot[plotstyle=line,showpoints=true,linecolor=gray,linewidth=0.7pt, dotstyle=|, dotangle=90, dotsize=0.15, linestyle=dotted]{%
    1 18.9
2 18.4
3 19.7
4 18.8
}

% Moyenne de la classe
\listplot[plotstyle=line,showpoints=true,linecolor=gray,linewidth=0.7pt, dotstyle=x, dotsize = 0.15, linestyle=dashed]{%
    1 12.4466666666667
2 13.1933333333333
3 10.82
4 11.5066666666667
}
\end{pspicture}
\pagebreak
\begin{tcolorbox}[enhanced,width=\textwidth,center upper,fontupper=\bfseries,drop shadow southwest,sharp corners]
{\sc \large Body} Timothée
\end{tcolorbox}
\medskip
\begin{tabularx}{\textwidth}{p{5cm}X}
	\alertbox{\faAward}{Note}{
		\begin{itemize}[leftmargin=0pt]
			\item[\textbullet] Note : \textbf{\large 18.8}
			\item[\textbullet] Rang : \textbf{1}
			\item[\textbullet] Traité : 97 \%
		\end{itemize}
	} &
	\alertbox{\faChartLine}{Statistiques des notes}{
        \psset{xunit=1cm, yunit=1cm,fillstyle=solid}
		\begin{pspicture}(0,-0.1)(16,1.45)
		   \savedata{\data}[11.6 14.5 18.8 8.4 8.2 7.6 14.1 8.1 9.6 14.2 17.6 14.9 18.6 6.4 0.0]
		   \rput{-90}(0,0.9){\psBoxplot[barwidth=1.1cm,yunit=0.5,fillcolor=gray,linewidth=1pt]{\data}}
		   \psaxes[yAxis=false,dx=1cm,Dx=2,labelsep=1pt,linecolor=gray,xlabelFontSize=\scriptstyle](0,0)(10.1,4)
		   \psdot[dotsize=8pt,dotstyle=diamond,linecolor=black,fillstyle=solid,fillcolor=white,linewidth=1pt](9.4,0.85)
           \psdot[dotsize=6pt,dotstyle=x,linecolor=black,linewidth=3pt](5.753333333333333,0.85)
		   \end{pspicture}
	} \\
    
\end{tabularx}\\
\begin{tabularx}{\textwidth}{X}
\alertbox{\faComment}{Commentaire}
{
	Excellent travail une fois de plus, bravo !
C’est avant tout le problème des dépassement de capacité (UB sur les entiers signés) qui justifient l’utilisation d’un type non signé.
Je n’ai pas trouvé trace de la question 7, un oubli ?
}
\end{tabularx}
\medskip
    \ding{113} \textbf{\sffamily{Résultats par thème}} \medskip \\
    \renewcommand{\arraystretch}{1.2}
    \begin{tabular}{|l|r|r|}
    \cline{2-3}
    \multicolumn{1}{l|}{} & \multicolumn{1}{|c|}{Points} & \multicolumn{1}{|c|}{Traitées} \\
    \hline
    {Base d'OCaml (fonctionnel)} & 77\% \;{\small (35/45)} & 80\% \;{\small (4/5)} \\ \hline {Complexité d'un algorithme} & 100\% \;{\small (30/30)} & 100\% \;{\small (4/4)} \\ \hline {Représentation des entiers} & 40\% \;{\small (02/5)} & 100\% \;{\small (1/1)} \\ \hline {Structure de données séquentielles} & 100\% \;{\small (45/45)} & 100\% \;{\small (6/6)} \\ \hline {Base de données et SQL} & 100\% \;{\small (60/60)} & 100\% \;{\small (8/8)} \\ \hline {Comprendre un algorithme} & 100\% \;{\small (20/20)} & 100\% \;{\small (4/4)} \\ \hline {Types structurés en C et pointeurs} & 100\% \;{\small (20/20)} & 100\% \;{\small (2/2)} \\ \hline {Table de hachage} & 90\% \;{\small (36/40)} & 100\% \;{\small (4/4)} \\ \hline {Programmation de base en C} & 100\% \;{\small (15/15)} & 100\% \;{\small (2/2)} \\ \hline \end{tabular} \\\\\medskip \\
    \ding{113} \textbf{\sffamily{Résultats par exercice}} \medskip \\
    \renewcommand{\arraystretch}{1.2}
    \begin{tabular}{|l|r|r|}
    \cline{2-3}
    \multicolumn{1}{l|}{} & \multicolumn{1}{|c|}{Points} & \multicolumn{1}{|c|}{Traitées} \\
    \hline
    Exercice {1} & 81\% \;{\small (45/55)} & 85\% \;{\small (6/7)} \\ \hline Exercice {2} & 100\% \;{\small (50/50)} & 100\% \;{\small (7/7)} \\ \hline Exercice {3} & 100\% \;{\small (60/60)} & 100\% \;{\small (8/8)} \\ \hline Exercice {4} & 87\% \;{\small (48/55)} & 100\% \;{\small (6/6)} \\ \hline Exercice {5} & 100\% \;{\small (60/60)} & 100\% \;{\small (8/8)} \\ \hline \end{tabular} \\\\
    \vspace{0.5cm}\\
    \ding{113} \textbf{\sffamily{Historique des notes}} \medskip \\
    \psset{xunit=1.4cm, yunit=0.2cm}
    \begin{pspicture}(-1,-1)(12,22)

% --- Axe X : numéros de devoir ---
\psaxes[
    Dx=1,
    Dy=2,
    Ox=0,
    Oy=0,
    labels=all,
    ticks=all,
](0,0)(0,0)(7,20)


\listplot[plotstyle=line,showpoints=true,linecolor=black,linewidth=0.7pt, dotstyle=diamond*, dotsize=0.2]{%
    1 18.9
2 18.4
3 19.2
4 18.8
}

% Minimum de la classe
\listplot[plotstyle=line,showpoints=true,linecolor=gray,linewidth=0.7pt, dotstyle=|, dotangle=90, dotsize=0.15, linestyle=dotted]{%
    1 5.8
2 4.6
3 0.9
4 0
}

% Maximum de la classe
\listplot[plotstyle=line,showpoints=true,linecolor=gray,linewidth=0.7pt, dotstyle=|, dotangle=90, dotsize=0.15, linestyle=dotted]{%
    1 18.9
2 18.4
3 19.7
4 18.8
}

% Moyenne de la classe
\listplot[plotstyle=line,showpoints=true,linecolor=gray,linewidth=0.7pt, dotstyle=x, dotsize = 0.15, linestyle=dashed]{%
    1 12.4466666666667
2 13.1933333333333
3 10.82
4 11.5066666666667
}
\end{pspicture}
\pagebreak
\begin{tcolorbox}[enhanced,width=\textwidth,center upper,fontupper=\bfseries,drop shadow southwest,sharp corners]
{\sc \large Boucher} Mathis
\end{tcolorbox}
\medskip
\begin{tabularx}{\textwidth}{p{5cm}X}
	\alertbox{\faAward}{Note}{
		\begin{itemize}[leftmargin=0pt]
			\item[\textbullet] Note : \textbf{\large 8.4}
			\item[\textbullet] Rang : \textbf{10}
			\item[\textbullet] Traité : 69 \%
		\end{itemize}
	} &
	\alertbox{\faChartLine}{Statistiques des notes}{
        \psset{xunit=1cm, yunit=1cm,fillstyle=solid}
		\begin{pspicture}(0,-0.1)(16,1.45)
		   \savedata{\data}[11.6 14.5 18.8 8.4 8.2 7.6 14.1 8.1 9.6 14.2 17.6 14.9 18.6 6.4 0.0]
		   \rput{-90}(0,0.9){\psBoxplot[barwidth=1.1cm,yunit=0.5,fillcolor=gray,linewidth=1pt]{\data}}
		   \psaxes[yAxis=false,dx=1cm,Dx=2,labelsep=1pt,linecolor=gray,xlabelFontSize=\scriptstyle](0,0)(10.1,4)
		   \psdot[dotsize=8pt,dotstyle=diamond,linecolor=black,fillstyle=solid,fillcolor=white,linewidth=1pt](4.2,0.85)
           \psdot[dotsize=6pt,dotstyle=x,linecolor=black,linewidth=3pt](5.753333333333333,0.85)
		   \end{pspicture}
	} \\
    
\end{tabularx}\\
\begin{tabularx}{\textwidth}{X}
\alertbox{\faComment}{Commentaire}
{
	L’ensemble est trop moyen, tu perds des points faciles sur le SQL, il faut t’entrainer davantage. Il faut aussi gagner en rapidité de façon à traiter plus de questions.
}
\end{tabularx}
\medskip
    \ding{113} \textbf{\sffamily{Résultats par thème}} \medskip \\
    \renewcommand{\arraystretch}{1.2}
    \begin{tabular}{|l|r|r|}
    \cline{2-3}
    \multicolumn{1}{l|}{} & \multicolumn{1}{|c|}{Points} & \multicolumn{1}{|c|}{Traitées} \\
    \hline
    {Base d'OCaml (fonctionnel)} & 4\% \;{\small (02/45)} & 20\% \;{\small (1/5)} \\ \hline {Complexité d'un algorithme} & 6\% \;{\small (02/30)} & 25\% \;{\small (1/4)} \\ \hline {Représentation des entiers} & 80\% \;{\small (04/5)} & 100\% \;{\small (1/1)} \\ \hline {Structure de données séquentielles} & 73\% \;{\small (33/45)} & 100\% \;{\small (6/6)} \\ \hline {Base de données et SQL} & 80\% \;{\small (48/60)} & 100\% \;{\small (8/8)} \\ \hline {Comprendre un algorithme} & 5\% \;{\small (01/20)} & 75\% \;{\small (3/4)} \\ \hline {Types structurés en C et pointeurs} & 60\% \;{\small (12/20)} & 100\% \;{\small (2/2)} \\ \hline {Table de hachage} & 5\% \;{\small (02/40)} & 25\% \;{\small (1/4)} \\ \hline {Programmation de base en C} & 93\% \;{\small (14/15)} & 100\% \;{\small (2/2)} \\ \hline \end{tabular} \\\\\medskip \\
    \ding{113} \textbf{\sffamily{Résultats par exercice}} \medskip \\
    \renewcommand{\arraystretch}{1.2}
    \begin{tabular}{|l|r|r|}
    \cline{2-3}
    \multicolumn{1}{l|}{} & \multicolumn{1}{|c|}{Points} & \multicolumn{1}{|c|}{Traitées} \\
    \hline
    Exercice {1} & 63\% \;{\small (35/55)} & 100\% \;{\small (7/7)} \\ \hline Exercice {2} & 38\% \;{\small (19/50)} & 85\% \;{\small (6/7)} \\ \hline Exercice {3} & 80\% \;{\small (48/60)} & 100\% \;{\small (8/8)} \\ \hline Exercice {4} & 29\% \;{\small (16/55)} & 50\% \;{\small (3/6)} \\ \hline Exercice {5} & 0\% \;{\small (00/60)} & 12\% \;{\small (1/8)} \\ \hline \end{tabular} \\\\
    \vspace{0.5cm}\\
    \ding{113} \textbf{\sffamily{Historique des notes}} \medskip \\
    \psset{xunit=1.4cm, yunit=0.2cm}
    \begin{pspicture}(-1,-1)(12,22)

% --- Axe X : numéros de devoir ---
\psaxes[
    Dx=1,
    Dy=2,
    Ox=0,
    Oy=0,
    labels=all,
    ticks=all,
](0,0)(0,0)(7,20)


\listplot[plotstyle=line,showpoints=true,linecolor=black,linewidth=0.7pt, dotstyle=diamond*, dotsize=0.2]{%
    1 8.4
2 11.8
3 8
4 8.4
}

% Minimum de la classe
\listplot[plotstyle=line,showpoints=true,linecolor=gray,linewidth=0.7pt, dotstyle=|, dotangle=90, dotsize=0.15, linestyle=dotted]{%
    1 5.8
2 4.6
3 0.9
4 0
}

% Maximum de la classe
\listplot[plotstyle=line,showpoints=true,linecolor=gray,linewidth=0.7pt, dotstyle=|, dotangle=90, dotsize=0.15, linestyle=dotted]{%
    1 18.9
2 18.4
3 19.7
4 18.8
}

% Moyenne de la classe
\listplot[plotstyle=line,showpoints=true,linecolor=gray,linewidth=0.7pt, dotstyle=x, dotsize = 0.15, linestyle=dashed]{%
    1 12.4466666666667
2 13.1933333333333
3 10.82
4 11.5066666666667
}
\end{pspicture}
\pagebreak
\begin{tcolorbox}[enhanced,width=\textwidth,center upper,fontupper=\bfseries,drop shadow southwest,sharp corners]
{\sc \large Chane-lock} Maxime
\end{tcolorbox}
\medskip
\begin{tabularx}{\textwidth}{p{5cm}X}
	\alertbox{\faAward}{Note}{
		\begin{itemize}[leftmargin=0pt]
			\item[\textbullet] Note : \textbf{\large 8.2}
			\item[\textbullet] Rang : \textbf{11}
			\item[\textbullet] Traité : 67 \%
		\end{itemize}
	} &
	\alertbox{\faChartLine}{Statistiques des notes}{
        \psset{xunit=1cm, yunit=1cm,fillstyle=solid}
		\begin{pspicture}(0,-0.1)(16,1.45)
		   \savedata{\data}[11.6 14.5 18.8 8.4 8.2 7.6 14.1 8.1 9.6 14.2 17.6 14.9 18.6 6.4 0.0]
		   \rput{-90}(0,0.9){\psBoxplot[barwidth=1.1cm,yunit=0.5,fillcolor=gray,linewidth=1pt]{\data}}
		   \psaxes[yAxis=false,dx=1cm,Dx=2,labelsep=1pt,linecolor=gray,xlabelFontSize=\scriptstyle](0,0)(10.1,4)
		   \psdot[dotsize=8pt,dotstyle=diamond,linecolor=black,fillstyle=solid,fillcolor=white,linewidth=1pt](4.1,0.85)
           \psdot[dotsize=6pt,dotstyle=x,linecolor=black,linewidth=3pt](5.753333333333333,0.85)
		   \end{pspicture}
	} \\
    
\end{tabularx}\\
\begin{tabularx}{\textwidth}{X}
\alertbox{\faComment}{Commentaire}
{
	Les questions de cours sont bien traitées c’est bien mais tu perds des points sur le SQL (il faut revoir l’utilisation de GROUP BY). Il faut aussi gagner en rapidité afin de traiter plus de questions.
}
\end{tabularx}
\medskip
    \ding{113} \textbf{\sffamily{Résultats par thème}} \medskip \\
    \renewcommand{\arraystretch}{1.2}
    \begin{tabular}{|l|r|r|}
    \cline{2-3}
    \multicolumn{1}{l|}{} & \multicolumn{1}{|c|}{Points} & \multicolumn{1}{|c|}{Traitées} \\
    \hline
    {Base d'OCaml (fonctionnel)} & 0\% \;{\small (00/45)} & 0\% \;{\small (0/5)} \\ \hline {Complexité d'un algorithme} & 16\% \;{\small (05/30)} & 25\% \;{\small (1/4)} \\ \hline {Représentation des entiers} & 80\% \;{\small (04/5)} & 100\% \;{\small (1/1)} \\ \hline {Structure de données séquentielles} & 100\% \;{\small (45/45)} & 100\% \;{\small (6/6)} \\ \hline {Base de données et SQL} & 61\% \;{\small (37/60)} & 100\% \;{\small (8/8)} \\ \hline {Comprendre un algorithme} & 45\% \;{\small (09/20)} & 75\% \;{\small (3/4)} \\ \hline {Types structurés en C et pointeurs} & 0\% \;{\small (00/20)} & 0\% \;{\small (0/2)} \\ \hline {Table de hachage} & 25\% \;{\small (10/40)} & 100\% \;{\small (4/4)} \\ \hline {Programmation de base en C} & 33\% \;{\small (05/15)} & 50\% \;{\small (1/2)} \\ \hline \end{tabular} \\\\\medskip \\
    \ding{113} \textbf{\sffamily{Résultats par exercice}} \medskip \\
    \renewcommand{\arraystretch}{1.2}
    \begin{tabular}{|l|r|r|}
    \cline{2-3}
    \multicolumn{1}{l|}{} & \multicolumn{1}{|c|}{Points} & \multicolumn{1}{|c|}{Traitées} \\
    \hline
    Exercice {1} & 81\% \;{\small (45/55)} & 85\% \;{\small (6/7)} \\ \hline Exercice {2} & 38\% \;{\small (19/50)} & 57\% \;{\small (4/7)} \\ \hline Exercice {3} & 61\% \;{\small (37/60)} & 100\% \;{\small (8/8)} \\ \hline Exercice {4} & 25\% \;{\small (14/55)} & 83\% \;{\small (5/6)} \\ \hline Exercice {5} & 0\% \;{\small (00/60)} & 12\% \;{\small (1/8)} \\ \hline \end{tabular} \\\\
    \vspace{0.5cm}\\
    \ding{113} \textbf{\sffamily{Historique des notes}} \medskip \\
    \psset{xunit=1.4cm, yunit=0.2cm}
    \begin{pspicture}(-1,-1)(12,22)

% --- Axe X : numéros de devoir ---
\psaxes[
    Dx=1,
    Dy=2,
    Ox=0,
    Oy=0,
    labels=all,
    ticks=all,
](0,0)(0,0)(7,20)


\listplot[plotstyle=line,showpoints=true,linecolor=black,linewidth=0.7pt, dotstyle=diamond*, dotsize=0.2]{%
    1 8.4
2 9
3 8.1
4 8.2
}

% Minimum de la classe
\listplot[plotstyle=line,showpoints=true,linecolor=gray,linewidth=0.7pt, dotstyle=|, dotangle=90, dotsize=0.15, linestyle=dotted]{%
    1 5.8
2 4.6
3 0.9
4 0
}

% Maximum de la classe
\listplot[plotstyle=line,showpoints=true,linecolor=gray,linewidth=0.7pt, dotstyle=|, dotangle=90, dotsize=0.15, linestyle=dotted]{%
    1 18.9
2 18.4
3 19.7
4 18.8
}

% Moyenne de la classe
\listplot[plotstyle=line,showpoints=true,linecolor=gray,linewidth=0.7pt, dotstyle=x, dotsize = 0.15, linestyle=dashed]{%
    1 12.4466666666667
2 13.1933333333333
3 10.82
4 11.5066666666667
}
\end{pspicture}
\pagebreak
\begin{tcolorbox}[enhanced,width=\textwidth,center upper,fontupper=\bfseries,drop shadow southwest,sharp corners]
{\sc \large Courounadin-Mouny} Maxence
\end{tcolorbox}
\medskip
\begin{tabularx}{\textwidth}{p{5cm}X}
	\alertbox{\faAward}{Note}{
		\begin{itemize}[leftmargin=0pt]
			\item[\textbullet] Note : \textbf{\large 7.6}
			\item[\textbullet] Rang : \textbf{13}
			\item[\textbullet] Traité : 83 \%
		\end{itemize}
	} &
	\alertbox{\faChartLine}{Statistiques des notes}{
        \psset{xunit=1cm, yunit=1cm,fillstyle=solid}
		\begin{pspicture}(0,-0.1)(16,1.45)
		   \savedata{\data}[11.6 14.5 18.8 8.4 8.2 7.6 14.1 8.1 9.6 14.2 17.6 14.9 18.6 6.4 0.0]
		   \rput{-90}(0,0.9){\psBoxplot[barwidth=1.1cm,yunit=0.5,fillcolor=gray,linewidth=1pt]{\data}}
		   \psaxes[yAxis=false,dx=1cm,Dx=2,labelsep=1pt,linecolor=gray,xlabelFontSize=\scriptstyle](0,0)(10.1,4)
		   \psdot[dotsize=8pt,dotstyle=diamond,linecolor=black,fillstyle=solid,fillcolor=white,linewidth=1pt](3.8,0.85)
           \psdot[dotsize=6pt,dotstyle=x,linecolor=black,linewidth=3pt](5.753333333333333,0.85)
		   \end{pspicture}
	} \\
    
\end{tabularx}\\
\begin{tabularx}{\textwidth}{X}
\alertbox{\faComment}{Commentaire}
{
	Il faut davantage maitriser le cours pour atteindre la moyenne. On peut parcourir un arbre binaire même si ce n’est pas un arbre binaire de recherche !!!
Je te conseille de revoir l’exercice 2 tout en codant et testant les fonctions demandées. Le SQL demande aussi a être retravaillé, tu as perdu des points faciles.
}
\end{tabularx}
\medskip
    \ding{113} \textbf{\sffamily{Résultats par thème}} \medskip \\
    \renewcommand{\arraystretch}{1.2}
    \begin{tabular}{|l|r|r|}
    \cline{2-3}
    \multicolumn{1}{l|}{} & \multicolumn{1}{|c|}{Points} & \multicolumn{1}{|c|}{Traitées} \\
    \hline
    {Base d'OCaml (fonctionnel)} & 40\% \;{\small (18/45)} & 60\% \;{\small (3/5)} \\ \hline {Complexité d'un algorithme} & 16\% \;{\small (05/30)} & 50\% \;{\small (2/4)} \\ \hline {Représentation des entiers} & 100\% \;{\small (05/5)} & 100\% \;{\small (1/1)} \\ \hline {Structure de données séquentielles} & 33\% \;{\small (15/45)} & 100\% \;{\small (6/6)} \\ \hline {Base de données et SQL} & 56\% \;{\small (34/60)} & 100\% \;{\small (8/8)} \\ \hline {Comprendre un algorithme} & 50\% \;{\small (10/20)} & 100\% \;{\small (4/4)} \\ \hline {Types structurés en C et pointeurs} & 70\% \;{\small (14/20)} & 100\% \;{\small (2/2)} \\ \hline {Table de hachage} & 0\% \;{\small (00/40)} & 50\% \;{\small (2/4)} \\ \hline {Programmation de base en C} & 33\% \;{\small (05/15)} & 100\% \;{\small (2/2)} \\ \hline \end{tabular} \\\\\medskip \\
    \ding{113} \textbf{\sffamily{Résultats par exercice}} \medskip \\
    \renewcommand{\arraystretch}{1.2}
    \begin{tabular}{|l|r|r|}
    \cline{2-3}
    \multicolumn{1}{l|}{} & \multicolumn{1}{|c|}{Points} & \multicolumn{1}{|c|}{Traitées} \\
    \hline
    Exercice {1} & 41\% \;{\small (23/55)} & 100\% \;{\small (7/7)} \\ \hline Exercice {2} & 68\% \;{\small (34/50)} & 100\% \;{\small (7/7)} \\ \hline Exercice {3} & 56\% \;{\small (34/60)} & 100\% \;{\small (8/8)} \\ \hline Exercice {4} & 9\% \;{\small (05/55)} & 66\% \;{\small (4/6)} \\ \hline Exercice {5} & 16\% \;{\small (10/60)} & 50\% \;{\small (4/8)} \\ \hline \end{tabular} \\\\
    \vspace{0.5cm}\\
    \ding{113} \textbf{\sffamily{Historique des notes}} \medskip \\
    \psset{xunit=1.4cm, yunit=0.2cm}
    \begin{pspicture}(-1,-1)(12,22)

% --- Axe X : numéros de devoir ---
\psaxes[
    Dx=1,
    Dy=2,
    Ox=0,
    Oy=0,
    labels=all,
    ticks=all,
](0,0)(0,0)(7,20)


\listplot[plotstyle=line,showpoints=true,linecolor=black,linewidth=0.7pt, dotstyle=diamond*, dotsize=0.2]{%
    1 10.9
2 11.1
3 4.8
4 7.6
}

% Minimum de la classe
\listplot[plotstyle=line,showpoints=true,linecolor=gray,linewidth=0.7pt, dotstyle=|, dotangle=90, dotsize=0.15, linestyle=dotted]{%
    1 5.8
2 4.6
3 0.9
4 0
}

% Maximum de la classe
\listplot[plotstyle=line,showpoints=true,linecolor=gray,linewidth=0.7pt, dotstyle=|, dotangle=90, dotsize=0.15, linestyle=dotted]{%
    1 18.9
2 18.4
3 19.7
4 18.8
}

% Moyenne de la classe
\listplot[plotstyle=line,showpoints=true,linecolor=gray,linewidth=0.7pt, dotstyle=x, dotsize = 0.15, linestyle=dashed]{%
    1 12.4466666666667
2 13.1933333333333
3 10.82
4 11.5066666666667
}
\end{pspicture}
\pagebreak
\begin{tcolorbox}[enhanced,width=\textwidth,center upper,fontupper=\bfseries,drop shadow southwest,sharp corners]
{\sc \large Dominguez} Raphaël
\end{tcolorbox}
\medskip
\begin{tabularx}{\textwidth}{p{5cm}X}
	\alertbox{\faAward}{Note}{
		\begin{itemize}[leftmargin=0pt]
			\item[\textbullet] Note : \textbf{\large 14.1}
			\item[\textbullet] Rang : \textbf{7}
			\item[\textbullet] Traité : 86 \%
		\end{itemize}
	} &
	\alertbox{\faChartLine}{Statistiques des notes}{
        \psset{xunit=1cm, yunit=1cm,fillstyle=solid}
		\begin{pspicture}(0,-0.1)(16,1.45)
		   \savedata{\data}[11.6 14.5 18.8 8.4 8.2 7.6 14.1 8.1 9.6 14.2 17.6 14.9 18.6 6.4 0.0]
		   \rput{-90}(0,0.9){\psBoxplot[barwidth=1.1cm,yunit=0.5,fillcolor=gray,linewidth=1pt]{\data}}
		   \psaxes[yAxis=false,dx=1cm,Dx=2,labelsep=1pt,linecolor=gray,xlabelFontSize=\scriptstyle](0,0)(10.1,4)
		   \psdot[dotsize=8pt,dotstyle=diamond,linecolor=black,fillstyle=solid,fillcolor=white,linewidth=1pt](7.05,0.85)
           \psdot[dotsize=6pt,dotstyle=x,linecolor=black,linewidth=3pt](5.753333333333333,0.85)
		   \end{pspicture}
	} \\
    
\end{tabularx}\\
\begin{tabularx}{\textwidth}{X}
\alertbox{\faComment}{Commentaire}
{
	Bon travail dans l’ensemble, tu sembles avoir fait des progrès en langage C, c’est bien. Il faut maintenant travailler le Ocaml et aussi le SQL. Les aspects théoriques sont très bien traités, tu pourrais avoir d’excellents résultats en pratiquant davantage la programmation.
}
\end{tabularx}
\medskip
    \ding{113} \textbf{\sffamily{Résultats par thème}} \medskip \\
    \renewcommand{\arraystretch}{1.2}
    \begin{tabular}{|l|r|r|}
    \cline{2-3}
    \multicolumn{1}{l|}{} & \multicolumn{1}{|c|}{Points} & \multicolumn{1}{|c|}{Traitées} \\
    \hline
    {Base d'OCaml (fonctionnel)} & 26\% \;{\small (12/45)} & 80\% \;{\small (4/5)} \\ \hline {Complexité d'un algorithme} & 83\% \;{\small (25/30)} & 75\% \;{\small (3/4)} \\ \hline {Représentation des entiers} & 100\% \;{\small (05/5)} & 100\% \;{\small (1/1)} \\ \hline {Structure de données séquentielles} & 100\% \;{\small (45/45)} & 100\% \;{\small (6/6)} \\ \hline {Base de données et SQL} & 65\% \;{\small (39/60)} & 87\% \;{\small (7/8)} \\ \hline {Comprendre un algorithme} & 90\% \;{\small (18/20)} & 100\% \;{\small (4/4)} \\ \hline {Types structurés en C et pointeurs} & 100\% \;{\small (20/20)} & 100\% \;{\small (2/2)} \\ \hline {Table de hachage} & 70\% \;{\small (28/40)} & 75\% \;{\small (3/4)} \\ \hline {Programmation de base en C} & 33\% \;{\small (05/15)} & 50\% \;{\small (1/2)} \\ \hline \end{tabular} \\\\\medskip \\
    \ding{113} \textbf{\sffamily{Résultats par exercice}} \medskip \\
    \renewcommand{\arraystretch}{1.2}
    \begin{tabular}{|l|r|r|}
    \cline{2-3}
    \multicolumn{1}{l|}{} & \multicolumn{1}{|c|}{Points} & \multicolumn{1}{|c|}{Traitées} \\
    \hline
    Exercice {1} & 85\% \;{\small (47/55)} & 100\% \;{\small (7/7)} \\ \hline Exercice {2} & 100\% \;{\small (50/50)} & 100\% \;{\small (7/7)} \\ \hline Exercice {3} & 65\% \;{\small (39/60)} & 87\% \;{\small (7/8)} \\ \hline Exercice {4} & 60\% \;{\small (33/55)} & 66\% \;{\small (4/6)} \\ \hline Exercice {5} & 46\% \;{\small (28/60)} & 75\% \;{\small (6/8)} \\ \hline \end{tabular} \\\\
    \vspace{0.5cm}\\
    \ding{113} \textbf{\sffamily{Historique des notes}} \medskip \\
    \psset{xunit=1.4cm, yunit=0.2cm}
    \begin{pspicture}(-1,-1)(12,22)

% --- Axe X : numéros de devoir ---
\psaxes[
    Dx=1,
    Dy=2,
    Ox=0,
    Oy=0,
    labels=all,
    ticks=all,
](0,0)(0,0)(7,20)


\listplot[plotstyle=line,showpoints=true,linecolor=black,linewidth=0.7pt, dotstyle=diamond*, dotsize=0.2]{%
    1 15.7
2 16.1
3 6.6
4 14.1
}

% Minimum de la classe
\listplot[plotstyle=line,showpoints=true,linecolor=gray,linewidth=0.7pt, dotstyle=|, dotangle=90, dotsize=0.15, linestyle=dotted]{%
    1 5.8
2 4.6
3 0.9
4 0
}

% Maximum de la classe
\listplot[plotstyle=line,showpoints=true,linecolor=gray,linewidth=0.7pt, dotstyle=|, dotangle=90, dotsize=0.15, linestyle=dotted]{%
    1 18.9
2 18.4
3 19.7
4 18.8
}

% Moyenne de la classe
\listplot[plotstyle=line,showpoints=true,linecolor=gray,linewidth=0.7pt, dotstyle=x, dotsize = 0.15, linestyle=dashed]{%
    1 12.4466666666667
2 13.1933333333333
3 10.82
4 11.5066666666667
}
\end{pspicture}
\pagebreak
\begin{tcolorbox}[enhanced,width=\textwidth,center upper,fontupper=\bfseries,drop shadow southwest,sharp corners]
{\sc \large Garbal} Alizée
\end{tcolorbox}
\medskip
\begin{tabularx}{\textwidth}{p{5cm}X}
	\alertbox{\faAward}{Note}{
		\begin{itemize}[leftmargin=0pt]
			\item[\textbullet] Note : \textbf{\large 8.1}
			\item[\textbullet] Rang : \textbf{12}
			\item[\textbullet] Traité : 83 \%
		\end{itemize}
	} &
	\alertbox{\faChartLine}{Statistiques des notes}{
        \psset{xunit=1cm, yunit=1cm,fillstyle=solid}
		\begin{pspicture}(0,-0.1)(16,1.45)
		   \savedata{\data}[11.6 14.5 18.8 8.4 8.2 7.6 14.1 8.1 9.6 14.2 17.6 14.9 18.6 6.4 0.0]
		   \rput{-90}(0,0.9){\psBoxplot[barwidth=1.1cm,yunit=0.5,fillcolor=gray,linewidth=1pt]{\data}}
		   \psaxes[yAxis=false,dx=1cm,Dx=2,labelsep=1pt,linecolor=gray,xlabelFontSize=\scriptstyle](0,0)(10.1,4)
		   \psdot[dotsize=8pt,dotstyle=diamond,linecolor=black,fillstyle=solid,fillcolor=white,linewidth=1pt](4.05,0.85)
           \psdot[dotsize=6pt,dotstyle=x,linecolor=black,linewidth=3pt](5.753333333333333,0.85)
		   \end{pspicture}
	} \\
    
\end{tabularx}\\
\begin{tabularx}{\textwidth}{X}
\alertbox{\faComment}{Commentaire}
{
	Des progrès très nets par rapport aux évaluations précédentes, il faut continuer sur cette voie.
}
\end{tabularx}
\medskip
    \ding{113} \textbf{\sffamily{Résultats par thème}} \medskip \\
    \renewcommand{\arraystretch}{1.2}
    \begin{tabular}{|l|r|r|}
    \cline{2-3}
    \multicolumn{1}{l|}{} & \multicolumn{1}{|c|}{Points} & \multicolumn{1}{|c|}{Traitées} \\
    \hline
    {Base d'OCaml (fonctionnel)} & 26\% \;{\small (12/45)} & 80\% \;{\small (4/5)} \\ \hline {Complexité d'un algorithme} & 60\% \;{\small (18/30)} & 100\% \;{\small (4/4)} \\ \hline {Représentation des entiers} & 100\% \;{\small (05/5)} & 100\% \;{\small (1/1)} \\ \hline {Structure de données séquentielles} & 73\% \;{\small (33/45)} & 100\% \;{\small (6/6)} \\ \hline {Base de données et SQL} & 63\% \;{\small (38/60)} & 100\% \;{\small (8/8)} \\ \hline {Comprendre un algorithme} & 40\% \;{\small (08/20)} & 75\% \;{\small (3/4)} \\ \hline {Types structurés en C et pointeurs} & 0\% \;{\small (00/20)} & 100\% \;{\small (2/2)} \\ \hline {Table de hachage} & 0\% \;{\small (00/40)} & 0\% \;{\small (0/4)} \\ \hline {Programmation de base en C} & 0\% \;{\small (00/15)} & 100\% \;{\small (2/2)} \\ \hline \end{tabular} \\\\\medskip \\
    \ding{113} \textbf{\sffamily{Résultats par exercice}} \medskip \\
    \renewcommand{\arraystretch}{1.2}
    \begin{tabular}{|l|r|r|}
    \cline{2-3}
    \multicolumn{1}{l|}{} & \multicolumn{1}{|c|}{Points} & \multicolumn{1}{|c|}{Traitées} \\
    \hline
    Exercice {1} & 78\% \;{\small (43/55)} & 100\% \;{\small (7/7)} \\ \hline Exercice {2} & 46\% \;{\small (23/50)} & 100\% \;{\small (7/7)} \\ \hline Exercice {3} & 63\% \;{\small (38/60)} & 100\% \;{\small (8/8)} \\ \hline Exercice {4} & 9\% \;{\small (05/55)} & 33\% \;{\small (2/6)} \\ \hline Exercice {5} & 8\% \;{\small (05/60)} & 75\% \;{\small (6/8)} \\ \hline \end{tabular} \\\\
    \vspace{0.5cm}\\
    \ding{113} \textbf{\sffamily{Historique des notes}} \medskip \\
    \psset{xunit=1.4cm, yunit=0.2cm}
    \begin{pspicture}(-1,-1)(12,22)

% --- Axe X : numéros de devoir ---
\psaxes[
    Dx=1,
    Dy=2,
    Ox=0,
    Oy=0,
    labels=all,
    ticks=all,
](0,0)(0,0)(7,20)


\listplot[plotstyle=line,showpoints=true,linecolor=black,linewidth=0.7pt, dotstyle=diamond*, dotsize=0.2]{%
    1 5.8
2 4.6
3 0.9
4 8.1
}

% Minimum de la classe
\listplot[plotstyle=line,showpoints=true,linecolor=gray,linewidth=0.7pt, dotstyle=|, dotangle=90, dotsize=0.15, linestyle=dotted]{%
    1 5.8
2 4.6
3 0.9
4 0
}

% Maximum de la classe
\listplot[plotstyle=line,showpoints=true,linecolor=gray,linewidth=0.7pt, dotstyle=|, dotangle=90, dotsize=0.15, linestyle=dotted]{%
    1 18.9
2 18.4
3 19.7
4 18.8
}

% Moyenne de la classe
\listplot[plotstyle=line,showpoints=true,linecolor=gray,linewidth=0.7pt, dotstyle=x, dotsize = 0.15, linestyle=dashed]{%
    1 12.4466666666667
2 13.1933333333333
3 10.82
4 11.5066666666667
}
\end{pspicture}
\pagebreak
\begin{tcolorbox}[enhanced,width=\textwidth,center upper,fontupper=\bfseries,drop shadow southwest,sharp corners]
{\sc \large Hoarau} Alessandro
\end{tcolorbox}
\medskip
\begin{tabularx}{\textwidth}{p{5cm}X}
	\alertbox{\faAward}{Note}{
		\begin{itemize}[leftmargin=0pt]
			\item[\textbullet] Note : \textbf{\large 9.6}
			\item[\textbullet] Rang : \textbf{9}
			\item[\textbullet] Traité : 69 \%
		\end{itemize}
	} &
	\alertbox{\faChartLine}{Statistiques des notes}{
        \psset{xunit=1cm, yunit=1cm,fillstyle=solid}
		\begin{pspicture}(0,-0.1)(16,1.45)
		   \savedata{\data}[11.6 14.5 18.8 8.4 8.2 7.6 14.1 8.1 9.6 14.2 17.6 14.9 18.6 6.4 0.0]
		   \rput{-90}(0,0.9){\psBoxplot[barwidth=1.1cm,yunit=0.5,fillcolor=gray,linewidth=1pt]{\data}}
		   \psaxes[yAxis=false,dx=1cm,Dx=2,labelsep=1pt,linecolor=gray,xlabelFontSize=\scriptstyle](0,0)(10.1,4)
		   \psdot[dotsize=8pt,dotstyle=diamond,linecolor=black,fillstyle=solid,fillcolor=white,linewidth=1pt](4.8,0.85)
           \psdot[dotsize=6pt,dotstyle=x,linecolor=black,linewidth=3pt](5.753333333333333,0.85)
		   \end{pspicture}
	} \\
    
\end{tabularx}\\
\begin{tabularx}{\textwidth}{X}
\alertbox{\faComment}{Commentaire}
{
	Ensemble correct, il faut travailler davantage ton cours. Tu dois aussi gagner en rapidité afin de traiter plus de questions.
}
\end{tabularx}
\medskip
    \ding{113} \textbf{\sffamily{Résultats par thème}} \medskip \\
    \renewcommand{\arraystretch}{1.2}
    \begin{tabular}{|l|r|r|}
    \cline{2-3}
    \multicolumn{1}{l|}{} & \multicolumn{1}{|c|}{Points} & \multicolumn{1}{|c|}{Traitées} \\
    \hline
    {Base d'OCaml (fonctionnel)} & 28\% \;{\small (13/45)} & 60\% \;{\small (3/5)} \\ \hline {Complexité d'un algorithme} & 13\% \;{\small (04/30)} & 25\% \;{\small (1/4)} \\ \hline {Représentation des entiers} & 0\% \;{\small (00/5)} & 100\% \;{\small (1/1)} \\ \hline {Structure de données séquentielles} & 66\% \;{\small (30/45)} & 83\% \;{\small (5/6)} \\ \hline {Base de données et SQL} & 78\% \;{\small (47/60)} & 100\% \;{\small (8/8)} \\ \hline {Comprendre un algorithme} & 60\% \;{\small (12/20)} & 75\% \;{\small (3/4)} \\ \hline {Types structurés en C et pointeurs} & 70\% \;{\small (14/20)} & 100\% \;{\small (2/2)} \\ \hline {Table de hachage} & 0\% \;{\small (00/40)} & 0\% \;{\small (0/4)} \\ \hline {Programmation de base en C} & 100\% \;{\small (15/15)} & 100\% \;{\small (2/2)} \\ \hline \end{tabular} \\\\\medskip \\
    \ding{113} \textbf{\sffamily{Résultats par exercice}} \medskip \\
    \renewcommand{\arraystretch}{1.2}
    \begin{tabular}{|l|r|r|}
    \cline{2-3}
    \multicolumn{1}{l|}{} & \multicolumn{1}{|c|}{Points} & \multicolumn{1}{|c|}{Traitées} \\
    \hline
    Exercice {1} & 65\% \;{\small (36/55)} & 85\% \;{\small (6/7)} \\ \hline Exercice {2} & 66\% \;{\small (33/50)} & 85\% \;{\small (6/7)} \\ \hline Exercice {3} & 78\% \;{\small (47/60)} & 100\% \;{\small (8/8)} \\ \hline Exercice {4} & 18\% \;{\small (10/55)} & 33\% \;{\small (2/6)} \\ \hline Exercice {5} & 15\% \;{\small (09/60)} & 37\% \;{\small (3/8)} \\ \hline \end{tabular} \\\\
    \vspace{0.5cm}\\
    \ding{113} \textbf{\sffamily{Historique des notes}} \medskip \\
    \psset{xunit=1.4cm, yunit=0.2cm}
    \begin{pspicture}(-1,-1)(12,22)

% --- Axe X : numéros de devoir ---
\psaxes[
    Dx=1,
    Dy=2,
    Ox=0,
    Oy=0,
    labels=all,
    ticks=all,
](0,0)(0,0)(7,20)


\listplot[plotstyle=line,showpoints=true,linecolor=black,linewidth=0.7pt, dotstyle=diamond*, dotsize=0.2]{%
    1 8
2 8.2
3 7.5
4 9.6
}

% Minimum de la classe
\listplot[plotstyle=line,showpoints=true,linecolor=gray,linewidth=0.7pt, dotstyle=|, dotangle=90, dotsize=0.15, linestyle=dotted]{%
    1 5.8
2 4.6
3 0.9
4 0
}

% Maximum de la classe
\listplot[plotstyle=line,showpoints=true,linecolor=gray,linewidth=0.7pt, dotstyle=|, dotangle=90, dotsize=0.15, linestyle=dotted]{%
    1 18.9
2 18.4
3 19.7
4 18.8
}

% Moyenne de la classe
\listplot[plotstyle=line,showpoints=true,linecolor=gray,linewidth=0.7pt, dotstyle=x, dotsize = 0.15, linestyle=dashed]{%
    1 12.4466666666667
2 13.1933333333333
3 10.82
4 11.5066666666667
}
\end{pspicture}
\pagebreak
\begin{tcolorbox}[enhanced,width=\textwidth,center upper,fontupper=\bfseries,drop shadow southwest,sharp corners]
{\sc \large Mahomed Issop} Jérémy
\end{tcolorbox}
\medskip
\begin{tabularx}{\textwidth}{p{5cm}X}
	\alertbox{\faAward}{Note}{
		\begin{itemize}[leftmargin=0pt]
			\item[\textbullet] Note : \textbf{\large 14.2}
			\item[\textbullet] Rang : \textbf{6}
			\item[\textbullet] Traité : 75 \%
		\end{itemize}
	} &
	\alertbox{\faChartLine}{Statistiques des notes}{
        \psset{xunit=1cm, yunit=1cm,fillstyle=solid}
		\begin{pspicture}(0,-0.1)(16,1.45)
		   \savedata{\data}[11.6 14.5 18.8 8.4 8.2 7.6 14.1 8.1 9.6 14.2 17.6 14.9 18.6 6.4 0.0]
		   \rput{-90}(0,0.9){\psBoxplot[barwidth=1.1cm,yunit=0.5,fillcolor=gray,linewidth=1pt]{\data}}
		   \psaxes[yAxis=false,dx=1cm,Dx=2,labelsep=1pt,linecolor=gray,xlabelFontSize=\scriptstyle](0,0)(10.1,4)
		   \psdot[dotsize=8pt,dotstyle=diamond,linecolor=black,fillstyle=solid,fillcolor=white,linewidth=1pt](7.1,0.85)
           \psdot[dotsize=6pt,dotstyle=x,linecolor=black,linewidth=3pt](5.753333333333333,0.85)
		   \end{pspicture}
	} \\
    
\end{tabularx}\\
\begin{tabularx}{\textwidth}{X}
\alertbox{\faComment}{Commentaire}
{
	Bon travail, c’est dommage que tu n’aies pas eu le temps de traiter l’exercice 5, il faut gagner en rapidité !
}
\end{tabularx}
\medskip
    \ding{113} \textbf{\sffamily{Résultats par thème}} \medskip \\
    \renewcommand{\arraystretch}{1.2}
    \begin{tabular}{|l|r|r|}
    \cline{2-3}
    \multicolumn{1}{l|}{} & \multicolumn{1}{|c|}{Points} & \multicolumn{1}{|c|}{Traitées} \\
    \hline
    {Base d'OCaml (fonctionnel)} & 22\% \;{\small (10/45)} & 20\% \;{\small (1/5)} \\ \hline {Complexité d'un algorithme} & 16\% \;{\small (05/30)} & 25\% \;{\small (1/4)} \\ \hline {Représentation des entiers} & 0\% \;{\small (00/5)} & 100\% \;{\small (1/1)} \\ \hline {Structure de données séquentielles} & 100\% \;{\small (45/45)} & 100\% \;{\small (6/6)} \\ \hline {Base de données et SQL} & 100\% \;{\small (60/60)} & 100\% \;{\small (8/8)} \\ \hline {Comprendre un algorithme} & 50\% \;{\small (10/20)} & 50\% \;{\small (2/4)} \\ \hline {Types structurés en C et pointeurs} & 80\% \;{\small (16/20)} & 100\% \;{\small (2/2)} \\ \hline {Table de hachage} & 95\% \;{\small (38/40)} & 100\% \;{\small (4/4)} \\ \hline {Programmation de base en C} & 100\% \;{\small (15/15)} & 100\% \;{\small (2/2)} \\ \hline \end{tabular} \\\\\medskip \\
    \ding{113} \textbf{\sffamily{Résultats par exercice}} \medskip \\
    \renewcommand{\arraystretch}{1.2}
    \begin{tabular}{|l|r|r|}
    \cline{2-3}
    \multicolumn{1}{l|}{} & \multicolumn{1}{|c|}{Points} & \multicolumn{1}{|c|}{Traitées} \\
    \hline
    Exercice {1} & 100\% \;{\small (55/55)} & 100\% \;{\small (7/7)} \\ \hline Exercice {2} & 72\% \;{\small (36/50)} & 85\% \;{\small (6/7)} \\ \hline Exercice {3} & 100\% \;{\small (60/60)} & 100\% \;{\small (8/8)} \\ \hline Exercice {4} & 87\% \;{\small (48/55)} & 100\% \;{\small (6/6)} \\ \hline Exercice {5} & 0\% \;{\small (00/60)} & 0\% \;{\small (0/8)} \\ \hline \end{tabular} \\\\
    \vspace{0.5cm}\\
    \ding{113} \textbf{\sffamily{Historique des notes}} \medskip \\
    \psset{xunit=1.4cm, yunit=0.2cm}
    \begin{pspicture}(-1,-1)(12,22)

% --- Axe X : numéros de devoir ---
\psaxes[
    Dx=1,
    Dy=2,
    Ox=0,
    Oy=0,
    labels=all,
    ticks=all,
](0,0)(0,0)(7,20)


\listplot[plotstyle=line,showpoints=true,linecolor=black,linewidth=0.7pt, dotstyle=diamond*, dotsize=0.2]{%
    1 13.5
2 15
3 14.8
4 14.2
}

% Minimum de la classe
\listplot[plotstyle=line,showpoints=true,linecolor=gray,linewidth=0.7pt, dotstyle=|, dotangle=90, dotsize=0.15, linestyle=dotted]{%
    1 5.8
2 4.6
3 0.9
4 0
}

% Maximum de la classe
\listplot[plotstyle=line,showpoints=true,linecolor=gray,linewidth=0.7pt, dotstyle=|, dotangle=90, dotsize=0.15, linestyle=dotted]{%
    1 18.9
2 18.4
3 19.7
4 18.8
}

% Moyenne de la classe
\listplot[plotstyle=line,showpoints=true,linecolor=gray,linewidth=0.7pt, dotstyle=x, dotsize = 0.15, linestyle=dashed]{%
    1 12.4466666666667
2 13.1933333333333
3 10.82
4 11.5066666666667
}
\end{pspicture}
\pagebreak
\begin{tcolorbox}[enhanced,width=\textwidth,center upper,fontupper=\bfseries,drop shadow southwest,sharp corners]
{\sc \large Mamodhoussen} Djavad
\end{tcolorbox}
\medskip
\begin{tabularx}{\textwidth}{p{5cm}X}
	\alertbox{\faAward}{Note}{
		\begin{itemize}[leftmargin=0pt]
			\item[\textbullet] Note : \textbf{\large 17.6}
			\item[\textbullet] Rang : \textbf{3}
			\item[\textbullet] Traité : 100 \%
		\end{itemize}
	} &
	\alertbox{\faChartLine}{Statistiques des notes}{
        \psset{xunit=1cm, yunit=1cm,fillstyle=solid}
		\begin{pspicture}(0,-0.1)(16,1.45)
		   \savedata{\data}[11.6 14.5 18.8 8.4 8.2 7.6 14.1 8.1 9.6 14.2 17.6 14.9 18.6 6.4 0.0]
		   \rput{-90}(0,0.9){\psBoxplot[barwidth=1.1cm,yunit=0.5,fillcolor=gray,linewidth=1pt]{\data}}
		   \psaxes[yAxis=false,dx=1cm,Dx=2,labelsep=1pt,linecolor=gray,xlabelFontSize=\scriptstyle](0,0)(10.1,4)
		   \psdot[dotsize=8pt,dotstyle=diamond,linecolor=black,fillstyle=solid,fillcolor=white,linewidth=1pt](8.8,0.85)
           \psdot[dotsize=6pt,dotstyle=x,linecolor=black,linewidth=3pt](5.753333333333333,0.85)
		   \end{pspicture}
	} \\
    
\end{tabularx}\\
\begin{tabularx}{\textwidth}{X}
\alertbox{\faComment}{Commentaire}
{
	Très bon travail, il faut continuer sur cette voie !
}
\end{tabularx}
\medskip
    \ding{113} \textbf{\sffamily{Résultats par thème}} \medskip \\
    \renewcommand{\arraystretch}{1.2}
    \begin{tabular}{|l|r|r|}
    \cline{2-3}
    \multicolumn{1}{l|}{} & \multicolumn{1}{|c|}{Points} & \multicolumn{1}{|c|}{Traitées} \\
    \hline
    {Base d'OCaml (fonctionnel)} & 71\% \;{\small (32/45)} & 100\% \;{\small (5/5)} \\ \hline {Complexité d'un algorithme} & 100\% \;{\small (30/30)} & 100\% \;{\small (4/4)} \\ \hline {Représentation des entiers} & 40\% \;{\small (02/5)} & 100\% \;{\small (1/1)} \\ \hline {Structure de données séquentielles} & 100\% \;{\small (45/45)} & 100\% \;{\small (6/6)} \\ \hline {Base de données et SQL} & 93\% \;{\small (56/60)} & 100\% \;{\small (8/8)} \\ \hline {Comprendre un algorithme} & 75\% \;{\small (15/20)} & 100\% \;{\small (4/4)} \\ \hline {Types structurés en C et pointeurs} & 90\% \;{\small (18/20)} & 100\% \;{\small (2/2)} \\ \hline {Table de hachage} & 95\% \;{\small (38/40)} & 100\% \;{\small (4/4)} \\ \hline {Programmation de base en C} & 66\% \;{\small (10/15)} & 100\% \;{\small (2/2)} \\ \hline \end{tabular} \\\\\medskip \\
    \ding{113} \textbf{\sffamily{Résultats par exercice}} \medskip \\
    \renewcommand{\arraystretch}{1.2}
    \begin{tabular}{|l|r|r|}
    \cline{2-3}
    \multicolumn{1}{l|}{} & \multicolumn{1}{|c|}{Points} & \multicolumn{1}{|c|}{Traitées} \\
    \hline
    Exercice {1} & 81\% \;{\small (45/55)} & 100\% \;{\small (7/7)} \\ \hline Exercice {2} & 86\% \;{\small (43/50)} & 100\% \;{\small (7/7)} \\ \hline Exercice {3} & 93\% \;{\small (56/60)} & 100\% \;{\small (8/8)} \\ \hline Exercice {4} & 90\% \;{\small (50/55)} & 100\% \;{\small (6/6)} \\ \hline Exercice {5} & 86\% \;{\small (52/60)} & 100\% \;{\small (8/8)} \\ \hline \end{tabular} \\\\
    \vspace{0.5cm}\\
    \ding{113} \textbf{\sffamily{Historique des notes}} \medskip \\
    \psset{xunit=1.4cm, yunit=0.2cm}
    \begin{pspicture}(-1,-1)(12,22)

% --- Axe X : numéros de devoir ---
\psaxes[
    Dx=1,
    Dy=2,
    Ox=0,
    Oy=0,
    labels=all,
    ticks=all,
](0,0)(0,0)(7,20)


\listplot[plotstyle=line,showpoints=true,linecolor=black,linewidth=0.7pt, dotstyle=diamond*, dotsize=0.2]{%
    1 17.8
2 18
3 16.4
4 17.6
}

% Minimum de la classe
\listplot[plotstyle=line,showpoints=true,linecolor=gray,linewidth=0.7pt, dotstyle=|, dotangle=90, dotsize=0.15, linestyle=dotted]{%
    1 5.8
2 4.6
3 0.9
4 0
}

% Maximum de la classe
\listplot[plotstyle=line,showpoints=true,linecolor=gray,linewidth=0.7pt, dotstyle=|, dotangle=90, dotsize=0.15, linestyle=dotted]{%
    1 18.9
2 18.4
3 19.7
4 18.8
}

% Moyenne de la classe
\listplot[plotstyle=line,showpoints=true,linecolor=gray,linewidth=0.7pt, dotstyle=x, dotsize = 0.15, linestyle=dashed]{%
    1 12.4466666666667
2 13.1933333333333
3 10.82
4 11.5066666666667
}
\end{pspicture}
\pagebreak
\begin{tcolorbox}[enhanced,width=\textwidth,center upper,fontupper=\bfseries,drop shadow southwest,sharp corners]
{\sc \large Morel} Lucas
\end{tcolorbox}
\medskip
\begin{tabularx}{\textwidth}{p{5cm}X}
	\alertbox{\faAward}{Note}{
		\begin{itemize}[leftmargin=0pt]
			\item[\textbullet] Note : \textbf{\large 14.9}
			\item[\textbullet] Rang : \textbf{4}
			\item[\textbullet] Traité : 92 \%
		\end{itemize}
	} &
	\alertbox{\faChartLine}{Statistiques des notes}{
        \psset{xunit=1cm, yunit=1cm,fillstyle=solid}
		\begin{pspicture}(0,-0.1)(16,1.45)
		   \savedata{\data}[11.6 14.5 18.8 8.4 8.2 7.6 14.1 8.1 9.6 14.2 17.6 14.9 18.6 6.4 0.0]
		   \rput{-90}(0,0.9){\psBoxplot[barwidth=1.1cm,yunit=0.5,fillcolor=gray,linewidth=1pt]{\data}}
		   \psaxes[yAxis=false,dx=1cm,Dx=2,labelsep=1pt,linecolor=gray,xlabelFontSize=\scriptstyle](0,0)(10.1,4)
		   \psdot[dotsize=8pt,dotstyle=diamond,linecolor=black,fillstyle=solid,fillcolor=white,linewidth=1pt](7.45,0.85)
           \psdot[dotsize=6pt,dotstyle=x,linecolor=black,linewidth=3pt](5.753333333333333,0.85)
		   \end{pspicture}
	} \\
    
\end{tabularx}\\
\begin{tabularx}{\textwidth}{X}
\alertbox{\faComment}{Commentaire}
{
	Très bon travail. Ne néglige pas le OCaml, par exemple, il faut revoir la représentation des arbres binaires en Ocaml (voir l’exercice de TP associé). Je te conseille aussi de retravailler l’exercice 4 (voir le corrigé en ligne si besoin). 
}
\end{tabularx}
\medskip
    \ding{113} \textbf{\sffamily{Résultats par thème}} \medskip \\
    \renewcommand{\arraystretch}{1.2}
    \begin{tabular}{|l|r|r|}
    \cline{2-3}
    \multicolumn{1}{l|}{} & \multicolumn{1}{|c|}{Points} & \multicolumn{1}{|c|}{Traitées} \\
    \hline
    {Base d'OCaml (fonctionnel)} & 66\% \;{\small (30/45)} & 100\% \;{\small (5/5)} \\ \hline {Complexité d'un algorithme} & 66\% \;{\small (20/30)} & 75\% \;{\small (3/4)} \\ \hline {Représentation des entiers} & 100\% \;{\small (05/5)} & 100\% \;{\small (1/1)} \\ \hline {Structure de données séquentielles} & 100\% \;{\small (45/45)} & 100\% \;{\small (6/6)} \\ \hline {Base de données et SQL} & 93\% \;{\small (56/60)} & 100\% \;{\small (8/8)} \\ \hline {Comprendre un algorithme} & 70\% \;{\small (14/20)} & 75\% \;{\small (3/4)} \\ \hline {Types structurés en C et pointeurs} & 80\% \;{\small (16/20)} & 100\% \;{\small (2/2)} \\ \hline {Table de hachage} & 45\% \;{\small (18/40)} & 100\% \;{\small (4/4)} \\ \hline {Programmation de base en C} & 33\% \;{\small (05/15)} & 50\% \;{\small (1/2)} \\ \hline \end{tabular} \\\\\medskip \\
    \ding{113} \textbf{\sffamily{Résultats par exercice}} \medskip \\
    \renewcommand{\arraystretch}{1.2}
    \begin{tabular}{|l|r|r|}
    \cline{2-3}
    \multicolumn{1}{l|}{} & \multicolumn{1}{|c|}{Points} & \multicolumn{1}{|c|}{Traitées} \\
    \hline
    Exercice {1} & 89\% \;{\small (49/55)} & 100\% \;{\small (7/7)} \\ \hline Exercice {2} & 72\% \;{\small (36/50)} & 85\% \;{\small (6/7)} \\ \hline Exercice {3} & 93\% \;{\small (56/60)} & 100\% \;{\small (8/8)} \\ \hline Exercice {4} & 41\% \;{\small (23/55)} & 83\% \;{\small (5/6)} \\ \hline Exercice {5} & 75\% \;{\small (45/60)} & 87\% \;{\small (7/8)} \\ \hline \end{tabular} \\\\
    \vspace{0.5cm}\\
    \ding{113} \textbf{\sffamily{Historique des notes}} \medskip \\
    \psset{xunit=1.4cm, yunit=0.2cm}
    \begin{pspicture}(-1,-1)(12,22)

% --- Axe X : numéros de devoir ---
\psaxes[
    Dx=1,
    Dy=2,
    Ox=0,
    Oy=0,
    labels=all,
    ticks=all,
](0,0)(0,0)(7,20)


\listplot[plotstyle=line,showpoints=true,linecolor=black,linewidth=0.7pt, dotstyle=diamond*, dotsize=0.2]{%
    1 16.7
2 15.5
3 13.6
4 14.9
}

% Minimum de la classe
\listplot[plotstyle=line,showpoints=true,linecolor=gray,linewidth=0.7pt, dotstyle=|, dotangle=90, dotsize=0.15, linestyle=dotted]{%
    1 5.8
2 4.6
3 0.9
4 0
}

% Maximum de la classe
\listplot[plotstyle=line,showpoints=true,linecolor=gray,linewidth=0.7pt, dotstyle=|, dotangle=90, dotsize=0.15, linestyle=dotted]{%
    1 18.9
2 18.4
3 19.7
4 18.8
}

% Moyenne de la classe
\listplot[plotstyle=line,showpoints=true,linecolor=gray,linewidth=0.7pt, dotstyle=x, dotsize = 0.15, linestyle=dashed]{%
    1 12.4466666666667
2 13.1933333333333
3 10.82
4 11.5066666666667
}
\end{pspicture}
\pagebreak
\begin{tcolorbox}[enhanced,width=\textwidth,center upper,fontupper=\bfseries,drop shadow southwest,sharp corners]
{\sc \large Randriamiarivola Korodo} Lionel
\end{tcolorbox}
\medskip
\begin{tabularx}{\textwidth}{p{5cm}X}
	\alertbox{\faAward}{Note}{
		\begin{itemize}[leftmargin=0pt]
			\item[\textbullet] Note : \textbf{\large 18.6}
			\item[\textbullet] Rang : \textbf{2}
			\item[\textbullet] Traité : 94 \%
		\end{itemize}
	} &
	\alertbox{\faChartLine}{Statistiques des notes}{
        \psset{xunit=1cm, yunit=1cm,fillstyle=solid}
		\begin{pspicture}(0,-0.1)(16,1.45)
		   \savedata{\data}[11.6 14.5 18.8 8.4 8.2 7.6 14.1 8.1 9.6 14.2 17.6 14.9 18.6 6.4 0.0]
		   \rput{-90}(0,0.9){\psBoxplot[barwidth=1.1cm,yunit=0.5,fillcolor=gray,linewidth=1pt]{\data}}
		   \psaxes[yAxis=false,dx=1cm,Dx=2,labelsep=1pt,linecolor=gray,xlabelFontSize=\scriptstyle](0,0)(10.1,4)
		   \psdot[dotsize=8pt,dotstyle=diamond,linecolor=black,fillstyle=solid,fillcolor=white,linewidth=1pt](9.3,0.85)
           \psdot[dotsize=6pt,dotstyle=x,linecolor=black,linewidth=3pt](5.753333333333333,0.85)
		   \end{pspicture}
	} \\
    
\end{tabularx}\\
\begin{tabularx}{\textwidth}{X}
\alertbox{\faComment}{Commentaire}
{
	Excellent travail, tes résultats sont vraiment prometteurs, continue sur cette voie !
}
\end{tabularx}
\medskip
    \ding{113} \textbf{\sffamily{Résultats par thème}} \medskip \\
    \renewcommand{\arraystretch}{1.2}
    \begin{tabular}{|l|r|r|}
    \cline{2-3}
    \multicolumn{1}{l|}{} & \multicolumn{1}{|c|}{Points} & \multicolumn{1}{|c|}{Traitées} \\
    \hline
    {Base d'OCaml (fonctionnel)} & 100\% \;{\small (45/45)} & 100\% \;{\small (5/5)} \\ \hline {Complexité d'un algorithme} & 66\% \;{\small (20/30)} & 75\% \;{\small (3/4)} \\ \hline {Représentation des entiers} & 100\% \;{\small (05/5)} & 100\% \;{\small (1/1)} \\ \hline {Structure de données séquentielles} & 100\% \;{\small (45/45)} & 100\% \;{\small (6/6)} \\ \hline {Base de données et SQL} & 100\% \;{\small (60/60)} & 100\% \;{\small (8/8)} \\ \hline {Comprendre un algorithme} & 100\% \;{\small (20/20)} & 100\% \;{\small (4/4)} \\ \hline {Types structurés en C et pointeurs} & 50\% \;{\small (10/20)} & 50\% \;{\small (1/2)} \\ \hline {Table de hachage} & 100\% \;{\small (40/40)} & 100\% \;{\small (4/4)} \\ \hline {Programmation de base en C} & 100\% \;{\small (15/15)} & 100\% \;{\small (2/2)} \\ \hline \end{tabular} \\\\\medskip \\
    \ding{113} \textbf{\sffamily{Résultats par exercice}} \medskip \\
    \renewcommand{\arraystretch}{1.2}
    \begin{tabular}{|l|r|r|}
    \cline{2-3}
    \multicolumn{1}{l|}{} & \multicolumn{1}{|c|}{Points} & \multicolumn{1}{|c|}{Traitées} \\
    \hline
    Exercice {1} & 100\% \;{\small (55/55)} & 100\% \;{\small (7/7)} \\ \hline Exercice {2} & 60\% \;{\small (30/50)} & 71\% \;{\small (5/7)} \\ \hline Exercice {3} & 100\% \;{\small (60/60)} & 100\% \;{\small (8/8)} \\ \hline Exercice {4} & 100\% \;{\small (55/55)} & 100\% \;{\small (6/6)} \\ \hline Exercice {5} & 100\% \;{\small (60/60)} & 100\% \;{\small (8/8)} \\ \hline \end{tabular} \\\\
    \vspace{0.5cm}\\
    \ding{113} \textbf{\sffamily{Historique des notes}} \medskip \\
    \psset{xunit=1.4cm, yunit=0.2cm}
    \begin{pspicture}(-1,-1)(12,22)

% --- Axe X : numéros de devoir ---
\psaxes[
    Dx=1,
    Dy=2,
    Ox=0,
    Oy=0,
    labels=all,
    ticks=all,
](0,0)(0,0)(7,20)


\listplot[plotstyle=line,showpoints=true,linecolor=black,linewidth=0.7pt, dotstyle=diamond*, dotsize=0.2]{%
    1 18.6
2 17.5
3 19.7
4 18.6
}

% Minimum de la classe
\listplot[plotstyle=line,showpoints=true,linecolor=gray,linewidth=0.7pt, dotstyle=|, dotangle=90, dotsize=0.15, linestyle=dotted]{%
    1 5.8
2 4.6
3 0.9
4 0
}

% Maximum de la classe
\listplot[plotstyle=line,showpoints=true,linecolor=gray,linewidth=0.7pt, dotstyle=|, dotangle=90, dotsize=0.15, linestyle=dotted]{%
    1 18.9
2 18.4
3 19.7
4 18.8
}

% Moyenne de la classe
\listplot[plotstyle=line,showpoints=true,linecolor=gray,linewidth=0.7pt, dotstyle=x, dotsize = 0.15, linestyle=dashed]{%
    1 12.4466666666667
2 13.1933333333333
3 10.82
4 11.5066666666667
}
\end{pspicture}
\pagebreak
\begin{tcolorbox}[enhanced,width=\textwidth,center upper,fontupper=\bfseries,drop shadow southwest,sharp corners]
{\sc \large Rasolofotsara} Ando
\end{tcolorbox}
\medskip
\begin{tabularx}{\textwidth}{p{5cm}X}
	\alertbox{\faAward}{Note}{
		\begin{itemize}[leftmargin=0pt]
			\item[\textbullet] Note : \textbf{\large 6.4}
			\item[\textbullet] Rang : \textbf{14}
			\item[\textbullet] Traité : 53 \%
		\end{itemize}
	} &
	\alertbox{\faChartLine}{Statistiques des notes}{
        \psset{xunit=1cm, yunit=1cm,fillstyle=solid}
		\begin{pspicture}(0,-0.1)(16,1.45)
		   \savedata{\data}[11.6 14.5 18.8 8.4 8.2 7.6 14.1 8.1 9.6 14.2 17.6 14.9 18.6 6.4 0.0]
		   \rput{-90}(0,0.9){\psBoxplot[barwidth=1.1cm,yunit=0.5,fillcolor=gray,linewidth=1pt]{\data}}
		   \psaxes[yAxis=false,dx=1cm,Dx=2,labelsep=1pt,linecolor=gray,xlabelFontSize=\scriptstyle](0,0)(10.1,4)
		   \psdot[dotsize=8pt,dotstyle=diamond,linecolor=black,fillstyle=solid,fillcolor=white,linewidth=1pt](3.2,0.85)
           \psdot[dotsize=6pt,dotstyle=x,linecolor=black,linewidth=3pt](5.753333333333333,0.85)
		   \end{pspicture}
	} \\
    
\end{tabularx}\\
\begin{tabularx}{\textwidth}{X}
\alertbox{\faComment}{Commentaire}
{
	Tu ne traites pas suffisamment de question, il faut gagner en rapidité et pour cela acquérir des automatismes. L’exercice 2 sur les piles est à revoir, je te conseille de le refaire à la maison en codant et testant les fonctions demandées.
}
\end{tabularx}
\medskip
    \ding{113} \textbf{\sffamily{Résultats par thème}} \medskip \\
    \renewcommand{\arraystretch}{1.2}
    \begin{tabular}{|l|r|r|}
    \cline{2-3}
    \multicolumn{1}{l|}{} & \multicolumn{1}{|c|}{Points} & \multicolumn{1}{|c|}{Traitées} \\
    \hline
    {Base d'OCaml (fonctionnel)} & 0\% \;{\small (00/45)} & 20\% \;{\small (1/5)} \\ \hline {Complexité d'un algorithme} & 0\% \;{\small (00/30)} & 0\% \;{\small (0/4)} \\ \hline {Représentation des entiers} & 0\% \;{\small (00/5)} & 100\% \;{\small (1/1)} \\ \hline {Structure de données séquentielles} & 66\% \;{\small (30/45)} & 83\% \;{\small (5/6)} \\ \hline {Base de données et SQL} & 65\% \;{\small (39/60)} & 87\% \;{\small (7/8)} \\ \hline {Comprendre un algorithme} & 40\% \;{\small (08/20)} & 50\% \;{\small (2/4)} \\ \hline {Types structurés en C et pointeurs} & 10\% \;{\small (02/20)} & 50\% \;{\small (1/2)} \\ \hline {Table de hachage} & 25\% \;{\small (10/40)} & 50\% \;{\small (2/4)} \\ \hline {Programmation de base en C} & 0\% \;{\small (00/15)} & 0\% \;{\small (0/2)} \\ \hline \end{tabular} \\\\\medskip \\
    \ding{113} \textbf{\sffamily{Résultats par exercice}} \medskip \\
    \renewcommand{\arraystretch}{1.2}
    \begin{tabular}{|l|r|r|}
    \cline{2-3}
    \multicolumn{1}{l|}{} & \multicolumn{1}{|c|}{Points} & \multicolumn{1}{|c|}{Traitées} \\
    \hline
    Exercice {1} & 54\% \;{\small (30/55)} & 85\% \;{\small (6/7)} \\ \hline Exercice {2} & 20\% \;{\small (10/50)} & 42\% \;{\small (3/7)} \\ \hline Exercice {3} & 65\% \;{\small (39/60)} & 87\% \;{\small (7/8)} \\ \hline Exercice {4} & 18\% \;{\small (10/55)} & 50\% \;{\small (3/6)} \\ \hline Exercice {5} & 0\% \;{\small (00/60)} & 0\% \;{\small (0/8)} \\ \hline \end{tabular} \\\\
    \vspace{0.5cm}\\
    \ding{113} \textbf{\sffamily{Historique des notes}} \medskip \\
    \psset{xunit=1.4cm, yunit=0.2cm}
    \begin{pspicture}(-1,-1)(12,22)

% --- Axe X : numéros de devoir ---
\psaxes[
    Dx=1,
    Dy=2,
    Ox=0,
    Oy=0,
    labels=all,
    ticks=all,
](0,0)(0,0)(7,20)


\listplot[plotstyle=line,showpoints=true,linecolor=black,linewidth=0.7pt, dotstyle=diamond*, dotsize=0.2]{%
    1 9.5
2 10
3 7
4 6.4
}

% Minimum de la classe
\listplot[plotstyle=line,showpoints=true,linecolor=gray,linewidth=0.7pt, dotstyle=|, dotangle=90, dotsize=0.15, linestyle=dotted]{%
    1 5.8
2 4.6
3 0.9
4 0
}

% Maximum de la classe
\listplot[plotstyle=line,showpoints=true,linecolor=gray,linewidth=0.7pt, dotstyle=|, dotangle=90, dotsize=0.15, linestyle=dotted]{%
    1 18.9
2 18.4
3 19.7
4 18.8
}

% Moyenne de la classe
\listplot[plotstyle=line,showpoints=true,linecolor=gray,linewidth=0.7pt, dotstyle=x, dotsize = 0.15, linestyle=dashed]{%
    1 12.4466666666667
2 13.1933333333333
3 10.82
4 11.5066666666667
}
\end{pspicture}
\pagebreak
\begin{tcolorbox}[enhanced,width=\textwidth,center upper,fontupper=\bfseries,drop shadow southwest,sharp corners]
{\sc \large Silotia} Donovan
\end{tcolorbox}
\medskip
\begin{tabularx}{\textwidth}{p{5cm}X}
	\alertbox{\faAward}{Note}{
		\begin{itemize}[leftmargin=0pt]
			\item[\textbullet] Note : \textbf{\large 0.0}
			\item[\textbullet] Rang : \textbf{15}
			\item[\textbullet] Traité : 0 \%
		\end{itemize}
	} &
	\alertbox{\faChartLine}{Statistiques des notes}{
        \psset{xunit=1cm, yunit=1cm,fillstyle=solid}
		\begin{pspicture}(0,-0.1)(16,1.45)
		   \savedata{\data}[11.6 14.5 18.8 8.4 8.2 7.6 14.1 8.1 9.6 14.2 17.6 14.9 18.6 6.4 0.0]
		   \rput{-90}(0,0.9){\psBoxplot[barwidth=1.1cm,yunit=0.5,fillcolor=gray,linewidth=1pt]{\data}}
		   \psaxes[yAxis=false,dx=1cm,Dx=2,labelsep=1pt,linecolor=gray,xlabelFontSize=\scriptstyle](0,0)(10.1,4)
		   \psdot[dotsize=8pt,dotstyle=diamond,linecolor=black,fillstyle=solid,fillcolor=white,linewidth=1pt](0.0,0.85)
           \psdot[dotsize=6pt,dotstyle=x,linecolor=black,linewidth=3pt](5.753333333333333,0.85)
		   \end{pspicture}
	} \\
    
\end{tabularx}\\
\begin{tabularx}{\textwidth}{X}
\alertbox{\faComment}{Commentaire}
{
	Absent
}
\end{tabularx}
\medskip
    \ding{113} \textbf{\sffamily{Résultats par thème}} \medskip \\
    \renewcommand{\arraystretch}{1.2}
    \begin{tabular}{|l|r|r|}
    \cline{2-3}
    \multicolumn{1}{l|}{} & \multicolumn{1}{|c|}{Points} & \multicolumn{1}{|c|}{Traitées} \\
    \hline
    {Base d'OCaml (fonctionnel)} & 0\% \;{\small (00/45)} & 0\% \;{\small (0/5)} \\ \hline {Complexité d'un algorithme} & 0\% \;{\small (00/30)} & 0\% \;{\small (0/4)} \\ \hline {Représentation des entiers} & 0\% \;{\small (00/5)} & 0\% \;{\small (0/1)} \\ \hline {Structure de données séquentielles} & 0\% \;{\small (00/45)} & 0\% \;{\small (0/6)} \\ \hline {Base de données et SQL} & 0\% \;{\small (00/60)} & 0\% \;{\small (0/8)} \\ \hline {Comprendre un algorithme} & 0\% \;{\small (00/20)} & 0\% \;{\small (0/4)} \\ \hline {Types structurés en C et pointeurs} & 0\% \;{\small (00/20)} & 0\% \;{\small (0/2)} \\ \hline {Table de hachage} & 0\% \;{\small (00/40)} & 0\% \;{\small (0/4)} \\ \hline {Programmation de base en C} & 0\% \;{\small (00/15)} & 0\% \;{\small (0/2)} \\ \hline \end{tabular} \\\\\medskip \\
    \ding{113} \textbf{\sffamily{Résultats par exercice}} \medskip \\
    \renewcommand{\arraystretch}{1.2}
    \begin{tabular}{|l|r|r|}
    \cline{2-3}
    \multicolumn{1}{l|}{} & \multicolumn{1}{|c|}{Points} & \multicolumn{1}{|c|}{Traitées} \\
    \hline
    Exercice {1} & 0\% \;{\small (00/55)} & 0\% \;{\small (0/7)} \\ \hline Exercice {2} & 0\% \;{\small (00/50)} & 0\% \;{\small (0/7)} \\ \hline Exercice {3} & 0\% \;{\small (00/60)} & 0\% \;{\small (0/8)} \\ \hline Exercice {4} & 0\% \;{\small (00/55)} & 0\% \;{\small (0/6)} \\ \hline Exercice {5} & 0\% \;{\small (00/60)} & 0\% \;{\small (0/8)} \\ \hline \end{tabular} \\\\
    \vspace{0.5cm}\\
    \ding{113} \textbf{\sffamily{Historique des notes}} \medskip \\
    \psset{xunit=1.4cm, yunit=0.2cm}
    \begin{pspicture}(-1,-1)(12,22)

% --- Axe X : numéros de devoir ---
\psaxes[
    Dx=1,
    Dy=2,
    Ox=0,
    Oy=0,
    labels=all,
    ticks=all,
](0,0)(0,0)(7,20)


\listplot[plotstyle=line,showpoints=true,linecolor=black,linewidth=0.7pt, dotstyle=diamond*, dotsize=0.2]{%
    1 11.6
2 15.4
3 9
4 0
}

% Minimum de la classe
\listplot[plotstyle=line,showpoints=true,linecolor=gray,linewidth=0.7pt, dotstyle=|, dotangle=90, dotsize=0.15, linestyle=dotted]{%
    1 5.8
2 4.6
3 0.9
4 0
}

% Maximum de la classe
\listplot[plotstyle=line,showpoints=true,linecolor=gray,linewidth=0.7pt, dotstyle=|, dotangle=90, dotsize=0.15, linestyle=dotted]{%
    1 18.9
2 18.4
3 19.7
4 18.8
}

% Moyenne de la classe
\listplot[plotstyle=line,showpoints=true,linecolor=gray,linewidth=0.7pt, dotstyle=x, dotsize = 0.15, linestyle=dashed]{%
    1 12.4466666666667
2 13.1933333333333
3 10.82
4 11.5066666666667
}
\end{pspicture}
\pagebreak\end{document}